\section{Open questions and future work}
\label{sec:future-work}

\subsection{Soundness analysis}

The implementation is complete and exercised end-to-end, but the
soundness of the lifted code in the verifier-side field
$F = \F_p[X]/\flift$ is not formally proved here. The structural
ingredients are in place --- the lifted code reduces mod $2$ to a
genuine $\GF(2^{16})$ Reed--Solomon code (\autoref{sec:lift}), which
is MDS with distance $N - n + 1$ --- but a complete analysis should:
\begin{itemize}
  \item Establish that the lifted code's effective minimum distance,
    as seen by the proximity test, matches (or closely tracks) the
    underlying $\GF(2^{16})$ RS distance, accounting for the
    structural lift's torsion.
  \item Carry the Zip$^+$ proximity argument through the
    $F = \F_p[X]/\flift$ setting. The existing argument uses a
    Schwartz--Zippel bound over a prime field; the analogue over the
    degree-$16$ extension $\F_{p^{16}}$ (valid when $\flift \bmod p$
    is irreducible, as it is for \code{P\_16\_DEFAULT}) is the natural
    target.
  \item Bound the soundness-error contribution of the radix-8
    evaluator specifically. The radix-8-with-$\GF$-lifted-Vandermonde
    code differs from a radix-2 $\Z$-lift in its higher $\Z$-bits
    (\autoref{sec:radix8}); the analysis should confirm the proximity
    argument is insensitive to that difference, since both reduce mod
    $2$ to the same $\GF(2^{16})$ code.
\end{itemize}

\subsection{Rate / opening-count tuning}

The shipped configurations use $\code{REP} = 4$ ($147$ column
openings) and $\code{REP} = 8$ ($100$ openings), matching the
protocol's IPRS opening table. The radix-8 evaluator constrains
$\log_2(\code{codeword\_len})$ to be a multiple of $3$, which couples
the admissible $(\code{row\_len}, \code{REP})$ pairs; e.g.\ a
bit-budget-matched $\code{REP} = 8$ run at $\code{num\_vars} = 10$
would need $\log_2(\code{codeword\_len}) = 13$, which the radix-8
constructor rejects. A radix-$2$/radix-$8$ \emph{mixed} evaluator
(one or two radix-2 stages plus radix-8 meta-stages) would lift this
divisibility constraint, at the cost of a slightly more intricate
kernel and a small reintroduction of radix-2 carry growth.

\subsection{Tighter codeword packing}

\code{BITS\_CW = 32} carries $\sim 6$ bits of headroom over the
empirical $\sim 2^{26}$ codeword-coefficient maximum at
$\code{codeword\_len} = 4096$. A workload-specific bound could push
\code{BITS\_CW} down toward $28$, shrinking the dominant proof
component a further $\sim 10\%$ --- at the cost of a non-power-of-two
field width and a tighter (workload-dependent) safety margin. The
coefficient magnitude grows with $\code{codeword\_len}$ (it approaches
$2^{63}$ at the maximal $\code{codeword\_len} = 2^{16}$, see
\autoref{sec:meas-coef}), so any tightened \code{BITS\_CW} must be
chosen per codeword length.

\subsection{Performance follow-ups}

\begin{itemize}
  \item \textbf{Within-FFT parallelism for small batches.} At
    $\code{batch} = 1$ the row-level rayon parallelism vanishes and
    only the per-meta-stage block loop remains; rayon's adaptive
    splitter may under-extract parallelism for small block counts.
  \item \textbf{SIMD matvec.} The $8 \times 8$ block matvec multiplies
    by $\{0,1\}$-coefficient lifted twiddles --- shift-and-add over
    \code{i64} coefficients --- which is a natural target for an
    explicit SIMD inner loop on hosts with wide vector registers.
  \item \textbf{Commit-side cost of radix-8.} The radix-8 form does
    more raw arithmetic than a radix-2 lift (\autoref{sec:opt-radix8}).
    The current trade favours proof size; a deployment that weights
    commit time more heavily could expose the radix as a tunable knob.
\end{itemize}

\subsection{Larger / other $\GF(2^k)$}

The implementation is fixed at $k = 16$, matching the protocol's
$1\times$-folded binary lane (\code{BinaryPoly<16>}). The algebra is
identical for other even $k$; only the precomputation changes ---
the irreducible $f$, the Cantor basis size, and the
\code{reduce\_mod\_ftilde} fold pattern. A $\GF(2^{32})$ instantiation
would serve the \emph{unfolded} binary lane directly (no protocol-side
folding), at the cost of larger codeword coefficients and a
$\GF(2^{32})$ evaluation domain.

\subsection{Code-survey reference}

For readers exploring the source, the key files are:

\begin{itemize}\small
  \item \code{zip-plus/src/code/binary\_add\_fft\_16.rs}: top-level,
    \code{BinaryAddFft16Code}, \code{LinearCode<Zt>} impl, smoke /
    roundtrip / measurement tests.
  \item \code{.../binary\_add\_fft\_16/basis.rs}:
    \code{Gf2\_16}, Cantor basis.
  \item \code{.../binary\_add\_fft\_16/params.rs}:
    \code{Config16}, precomputed basis / subspace polynomials.
  \item \code{.../binary\_add\_fft\_16/ring\_ops.rs}:
    \code{FftRingElement16}, lifted-ring arithmetic.
  \item \code{.../binary\_add\_fft\_16/radix8.rs}:
    \code{Radix8FftParams16}, the radix-8 kernel.
  \item \code{.../binary\_add\_fft\_16/ext\_field.rs}:
    \code{Gf2\_16Ext<P>} ($F = \F_p[X]/\flift$), \code{Bit4Poly16},
    \code{P\_16\_DEFAULT}, the irreducibility test.
  \item \code{.../binary\_add\_fft\_16/packed\_cw.rs}:
    \code{PackedI64Poly16<BITS>}, the bit-packed wire encoding.
  \item \code{zip-plus/benches/zip\_plus\_benches.rs}: the bench
    harness --- RAA / IPRS references plus the D16 R4 / R8 polychal
    groups.
\end{itemize}
